% GENERATED FILE. Edit canonical lessons, then run npm run generate:books.
% canonical-source-hash: fnv1a64:594f71aea92d166f
% canonical-lessons: GE-C516-buehne, GE-C516-proben, GE-C516-lampenfieber, GE-C516-auftreten, GE-C516-auftritt

\chapter{Auf der Bühne}
\label{ch:ge-auf-der-buhne}
\hlchaptermodality{516}

\begin{chapteropening}
\textbf{By the end of this chapter:} \emph{I can talk about performing on stage.}

Complete the last lesson of chapter 516: I can talk about performing on stage.
\end{chapteropening}

\section[die Bühne]{die Bühne — stage}

\label{lesson:GE-C516-buehne}



\begin{quote}
Before the new one: say the German for to conduct, then the German for to compose.
\end{quote}

\subsection*{You'll want to know: die Bühne}
\textbf{die Bühne} — \textquotedblleft{}stage\textquotedblright{}.

\begin{quote}
Die Band kommt um neun Uhr auf die Bühne.
\end{quote}

\begin{quote}
\emph{The band comes on stage at nine.}
\end{quote}

Plural: \emph{die Bühnen}.

\begin{cousinweb}
\emph{Bühne} is inherited from Middle High German \emph{büne} \textquotedblleft{}wooden stage, platform, ceiling\textquotedblright{} (EtymDB).
\end{cousinweb}

\begin{grammarlens}[title={What you've built}]
One more word for telling, explaining and describing at B1.
\end{grammarlens}

\subsection*{Your turn}
\begin{itemize}
\raggedright
  \item \emph{Say it:} \emph{die Bühne}
  \item \emph{Say it:} \emph{die Bühne}, once more
  \item \emph{Say it:} a sentence of your own with \emph{die Bühne}
  \item \emph{Recall:} say the German for to conduct, then the German for to compose, then say \emph{die Bühne} again
\end{itemize}

\begin{tcolorbox}[breakable,colback=teal!4,colframe=teal!35!black,title={Before you move on}]
What does die Bühne mean? (\textquotedblleft{}Stage\textquotedblright{}.) Say it once more.
\end{tcolorbox}
\section[proben]{proben — to rehearse}

\label{lesson:GE-C516-proben}



\begin{quote}
Before the new one: say the German for to compose, then the German for stage.
\end{quote}

\subsection*{You'll want to know: proben}
\textbf{proben} — \textquotedblleft{}to rehearse\textquotedblright{}.

\begin{quote}
Die Band probt jeden Freitag im Keller.
\end{quote}

\begin{quote}
\emph{The band rehearses every Friday in the cellar.}
\end{quote}

Regular: \emph{hat geprobt}. The noun \emph{die Probe} — rehearsal; also a test or a sample.

\begin{cousinweb}
\emph{proben} goes back to Latin \emph{probāre} \textquotedblleft{}to test, to try\textquotedblright{} (EtymDB).
\end{cousinweb}

\begin{grammarlens}[title={What you've built}]
One more word for telling, explaining and describing at B1.
\end{grammarlens}

\subsection*{Your turn}
\begin{itemize}
\raggedright
  \item \emph{Say it:} \emph{proben}
  \item \emph{Say it:} \emph{proben}, once more
  \item \emph{Say it:} a sentence of your own with \emph{proben}
  \item \emph{Recall:} say the German for to compose, then the German for stage, then say \emph{proben} again
\end{itemize}

\begin{tcolorbox}[breakable,colback=teal!4,colframe=teal!35!black,title={Before you move on}]
What does proben mean? (\textquotedblleft{}To rehearse\textquotedblright{}.) Say it once more.
\end{tcolorbox}
\section[das Lampenfieber]{das Lampenfieber — stage fright}

\label{lesson:GE-C516-lampenfieber}



\begin{quote}
Before the new one: say the German for stage, then the German for to rehearse.
\end{quote}

\subsection*{You'll want to know: das Lampenfieber}
\textbf{das Lampenfieber} — \textquotedblleft{}stage fright\textquotedblright{}.

\begin{quote}
Vor dem Konzert hatte die Sängerin Lampenfieber.
\end{quote}

\begin{quote}
\emph{Before the concert the singer had stage fright.}
\end{quote}

No plural. Literally lamp fever, after the stage lights.

\begin{cousinweb}
\emph{Lampenfieber} is built on \emph{Lampe}, from Latin \emph{lampas} \textquotedblleft{}torch, lamp\textquotedblright{}, and \emph{Fieber}, from Latin \emph{febris} \textquotedblleft{}fever\textquotedblright{} (EtymDB).
\end{cousinweb}

\begin{grammarlens}[title={What you've built}]
One more word for telling, explaining and describing at B1.
\end{grammarlens}

\subsection*{Your turn}
\begin{itemize}
\raggedright
  \item \emph{Say it:} \emph{das Lampenfieber}
  \item \emph{Say it:} \emph{das Lampenfieber}, once more
  \item \emph{Say it:} a sentence of your own with \emph{das Lampenfieber}
  \item \emph{Recall:} say the German for stage, then the German for to rehearse, then say \emph{das Lampenfieber} again
\end{itemize}

\begin{tcolorbox}[breakable,colback=teal!4,colframe=teal!35!black,title={Before you move on}]
What does das Lampenfieber mean? (\textquotedblleft{}Stage fright\textquotedblright{}.) Say it once more.
\end{tcolorbox}
\section[auftreten]{auftreten — to perform, to appear}

\label{lesson:GE-C516-auftreten}



\begin{quote}
Before the new one: say the German for to rehearse, then the German for stage fright.
\end{quote}

\subsection*{You'll want to know: auftreten}
\textbf{auftreten} — \textquotedblleft{}to perform, to appear\textquotedblright{}.

\begin{quote}
Die Band tritt heute Abend im Stadtpark auf.
\end{quote}

\begin{quote}
\emph{The band is performing in the city park tonight.}
\end{quote}

Separable, Perfekt with \emph{sein}: \emph{Sie ist oft im Fernsehen aufgetreten.} Also of problems: \emph{Es sind Probleme aufgetreten.}

\begin{cousinweb}
\emph{auftreten} is built on \emph{auf} \textquotedblleft{}up, on\textquotedblright{} and \emph{treten} \textquotedblleft{}to step\textquotedblright{}, inherited from Old High German \emph{tretan}, from Proto-Germanic \emph{tredaną} (EtymDB): to step up on stage.
\end{cousinweb}

\begin{grammarlens}[title={What you've built}]
One more word for telling, explaining and describing at B1.
\end{grammarlens}

\subsection*{Your turn}
\begin{itemize}
\raggedright
  \item \emph{Say it:} \emph{auftreten}
  \item \emph{Say it:} \emph{auftreten}, once more
  \item \emph{Say it:} a sentence of your own with \emph{auftreten}
  \item \emph{Recall:} say the German for to rehearse, then the German for stage fright, then say \emph{auftreten} again
\end{itemize}

\begin{tcolorbox}[breakable,colback=teal!4,colframe=teal!35!black,title={Before you move on}]
What does auftreten mean? (\textquotedblleft{}To perform, to appear\textquotedblright{}.) Say it once more.
\end{tcolorbox}
\section[der Auftritt]{der Auftritt — appearance, performance}

\label{lesson:GE-C516-auftritt}



\begin{quote}
Before the new one: say the German for stage fright, then the German for to perform.
\end{quote}

\subsection*{You'll want to know: der Auftritt}
\textbf{der Auftritt} — \textquotedblleft{}appearance, performance\textquotedblright{}.

\begin{quote}
Nach dem Auftritt waren alle Musiker erleichtert.
\end{quote}

\begin{quote}
\emph{After the performance all the musicians were relieved.}
\end{quote}

Plural: \emph{die Auftritte}.

\begin{cousinweb}
\emph{Auftritt} is built on \emph{auf} \textquotedblleft{}up, on\textquotedblright{} and \emph{Tritt} \textquotedblleft{}step\textquotedblright{}, from \emph{treten} \textquotedblleft{}to step\textquotedblright{}; EtymDB lists the German word without a source.
\end{cousinweb}

\begin{grammarlens}[title={What you've built}]
One more word for telling, explaining and describing at B1.
\end{grammarlens}

\subsection*{Your turn}
\begin{itemize}
\raggedright
  \item \emph{Say it:} \emph{der Auftritt}
  \item \emph{Say it:} \emph{der Auftritt}, once more
  \item \emph{Say it:} a sentence of your own with \emph{der Auftritt}
  \item \emph{Recall:} say the German for stage fright, then the German for to perform, then say \emph{der Auftritt} again
\end{itemize}

\begin{tcolorbox}[breakable,colback=teal!4,colframe=teal!35!black,title={Before you move on}]
What does der Auftritt mean? (\textquotedblleft{}Appearance, performance\textquotedblright{}.) Say it once more.
\end{tcolorbox}
